\begin{table*}[t]\centering\footnotesize
\caption{Joint evidence in the 31 Type I/II mitigation records. Required inapplicable items exclude a record from that query. The two Type III AHB records remain in the historical index but are excluded here.}\label{tab:r35-positive}
\begin{tabularx}{\textwidth}{@{}p{.19\textwidth}p{.38\textwidth}Y@{}}\toprule
Requested evidence & Bound record IDs & Interpretation at the source configuration\\\midrule
Inference + utility (3/20) & CACHEPRUNE-selective; CACHEPRUNE-context; SPARLEAK-replacement & Component quality; the two CachePrune targets share one package.\\
Inference + cost (5/21) & SPEC-padding-REST; DANGOL-capture; DANGOL-livebinding; both CACHEPRUNE records & Payload, transformed-capture, loopback and serving costs remain distinct.\\
Adaptation + inference (2/15) & UI-delay-offline; DANGOL-capture & Matched clean/protected fitting on a defended test.\\
Inference + utility + cost (2/19) & CACHEPRUNE-selective; CACHEPRUNE-context & One family and default configuration, two target outcomes.\\
All five items (0/13) & None in this extraction & No complete binding for this scoped query.\\\bottomrule
\end{tabularx}
\par\smallskip\raggedright SPEC reports a payload-size ratio at the evaluated configuration; its placement is only partially bound. This is not a measured wire-byte or end-to-end cost. Requiring explicit placement or a wire-cost object removes SPEC and leaves 4/21; the four retained records still use distinct cost objects.\end{table*}
